\documentclass[11pt,a4paper]{article}
\usepackage[T1]{fontenc}
\usepackage[spanish,es-noshorthands]{babel}
\usepackage[defaultfam]{montserrat}
\usepackage[a4paper,hmargin=23mm,top=24mm,bottom=22mm]{geometry}
\usepackage{xcolor}
\usepackage{microtype}
\usepackage{booktabs}
\usepackage[unicode,hidelinks]{hyperref}

\definecolor{brailleazul}{HTML}{016D84}
\definecolor{braillecereza}{HTML}{A3195B}
\definecolor{linea}{HTML}{CEDDDF}
\setlength{\parindent}{0pt}
\setlength{\parskip}{0pt}
\pagestyle{plain}

\newcommand{\codigo}[1]{{\fontfamily{\ttdefault}\fontseries{m}\selectfont\detokenize{#1}}}
\newcommand{\titulo}[1]{%
  {\small IES Luis Braille \hfill Bases de Datos · 1.º DAW}\par
  \vspace{5mm}
  {\fontsize{25}{30}\selectfont\bfseries\color{brailleazul}#1\par}
  \vspace{4mm}
  {\color{brailleazul}\hrule height 0.7pt}
  \vspace{6mm}
}
\newcommand{\apartado}[1]{%
  \par\addvspace{4mm}{\bfseries\color{brailleazul}#1}\par\vspace{1mm}%
}
\newenvironment{consulta}[2]{%
  \par\addvspace{6mm}%
  \pdfbookmark[0]{Consulta #1. #2}{consulta-#1}%
  \begin{minipage}{\linewidth}
    {\color{linea}\hrule height 0.5pt}\vspace{3mm}
    {\fontsize{16}{20}\selectfont\bfseries
      {\color{braillecereza}#1}\hspace{0.75em}{\color{brailleazul}#2}\par}%
    \vspace{2mm}%
}{%
  \end{minipage}\par
}
\newcommand{\salida}[2]{%
  \par\vspace{2mm}%
  {\small
    \hangindent=21mm \hangafter=1
    \makebox[21mm][l]{\bfseries Columnas}#1\par
    \vspace{0.7mm}%
    \hangindent=21mm \hangafter=1
    \makebox[21mm][l]{\bfseries Orden}#2\par
  }%
}

% Cada TXT contiene la cabecera y todas las filas devueltas por MySQL.
\ExplSyntaxOn
\ior_new:N \l_sakila_archivo_ior
\tl_new:N \l_sakila_linea_tl
\seq_new:N \l_sakila_campos_seq
\tl_new:N \l_sakila_cabecera_tl
\tl_new:N \l_sakila_filas_tl
\int_new:N \l_sakila_total_int
\int_const:Nn \c_sakila_limite_int {10}
\cs_generate_variant:Nn \seq_set_split:Nnn { Nxn, NxV }

\cs_new_protected:Npn \sakila_leer_resultados:n #1
  {
    \ior_open:Nn \l_sakila_archivo_ior { Resultados_06_Sakila/consulta-#1.txt }
    \ior_str_get:NN \l_sakila_archivo_ior \l_sakila_linea_tl
    \seq_set_split:NxV \l_sakila_campos_seq { \char_generate:nn {9} {12} } \l_sakila_linea_tl
    \seq_map_indexed_inline:Nn \l_sakila_campos_seq
      {
        \int_compare:nNnT {##1} > {1} { \tl_put_right:Nn \l_sakila_cabecera_tl { & } }
        \tl_put_right:Nn \l_sakila_cabecera_tl { \codigo{##2} }
      }

    \ior_str_map_inline:Nn \l_sakila_archivo_ior
      {
        \int_incr:N \l_sakila_total_int
        \int_compare:nNnF { \l_sakila_total_int } > { \c_sakila_limite_int }
          {
            \seq_set_split:Nxn \l_sakila_campos_seq { \char_generate:nn {9} {12} } { ##1 }
            \tl_put_right:Nx \l_sakila_filas_tl { \seq_use:Nn \l_sakila_campos_seq { & } \exp_not:N \\ }
          }
      }
    \ior_close:N \l_sakila_archivo_ior
  }

\NewDocumentCommand{\resultados}{mm}
  {
    \group_begin:
    \sakila_leer_resultados:n {#1}
    \par\addvspace{6mm}
    \small
    {\bfseries\color{brailleazul} Resultados~esperados}
    \enspace
    (primeras~
      \int_min:nn { \c_sakila_limite_int } { \l_sakila_total_int }~
      de~\int_use:N \l_sakila_total_int\ filas)
    \par\vspace{1mm}
    \renewcommand{\arraystretch}{1.12}
    \begin{tabular*}{\linewidth}{@{\extracolsep{\fill}}#2@{}}
      \toprule
      \l_sakila_cabecera_tl \\
      \midrule
      \l_sakila_filas_tl
      \bottomrule
    \end{tabular*}
    \par\vspace{2mm}
    \textbf{Resultado~completo}\enspace\codigo{consulta-#1.txt}\par
    \group_end:
  }
\ExplSyntaxOff

\begin{document}
\raggedright
\titulo{Consultas con JOIN en Sakila}

\apartado{Entrega}
Entrega las siete consultas \codigo{SELECT} en un único archivo
\codigo{practica_sakila.sql} y numéralas del 21 al 27 con comentarios.
Usa el tipo de \codigo{JOIN} indicado y respeta los nombres y el orden
de las columnas y de las filas.\par\smallskip
Usa MySQL 8 con la base de datos \codigo{sakila} y compara tu salida
con la muestra que aparece debajo de cada consulta. Los archivos de
\codigo{Resultados_06_Sakila} contienen todas las filas.

\apartado{Qué vas a practicar}
Composiciones externas, filtrado y ordenación de filas, funciones de agregado,
agrupación y filtrado de grupos. Usa \codigo{ON} para limitar las coincidencias,
\codigo{WHERE} para filtrar antes de agrupar y \codigo{HAVING} para filtrar los grupos.

\begin{consulta}{21}{Copias en la tienda 1}
Con \codigo{LEFT JOIN}, muestra las películas \codigo{PG} o \codigo{PG-13}
que tienen menos de tres copias en la tienda 1. Incluye las que tienen cero
copias en esa tienda, aunque tengan copias en la otra.
\salida{\codigo{film_id}, \codigo{title}, \codigo{rating}, \codigo{num_ejemplares}.}{Menor número de copias primero. En caso de empate, por título de A a Z.}
\end{consulta}
\resultados{21}{rlrr}
\clearpage

\begin{consulta}{22}{Películas sin copias en la tienda 2}
Con \codigo{RIGHT JOIN}, muestra las películas que duran al menos 120 minutos
y tienen cero copias en la tienda 2, aunque tengan copias en la otra.
\salida{\codigo{film_id}, \codigo{title}, \codigo{length}, \codigo{num_ejemplares}.}{Por título de A a Z.}
\end{consulta}
\resultados{22}{rlrr}
\clearpage

\begin{consulta}{23}{Clientes de la tienda 1 con pocos alquileres}
Con \codigo{LEFT JOIN}, muestra los clientes activos asignados a la tienda 1
que hicieron menos de tres alquileres en mayo de 2005. Incluye a quienes hicieron
cero. Cuenta los alquileres del cliente sin filtrar por la tienda de la copia.
\salida{\codigo{customer_id}, \codigo{first_name}, \codigo{last_name}, \codigo{num_alquileres}.}{De menos a más alquileres. Si empatan, por \codigo{customer_id} ascendente.}
\end{consulta}
\resultados{23}{rllr}
\clearpage

\begin{consulta}{24}{Pagos de clientes de la tienda 2}
Con \codigo{RIGHT JOIN}, muestra los clientes activos asignados a la tienda 2
con al menos cinco pagos individuales de importe igual o superior a 6,00.
Cuenta y suma solo los pagos que cumplen ese importe mínimo.
\salida{\codigo{customer_id}, \codigo{first_name}, \codigo{last_name}, \codigo{num_pagos}, \codigo{importe_total}.}{De más a menos pagos. Si empatan, por \codigo{customer_id} ascendente.}
\end{consulta}
\resultados{24}{rllrr}
\clearpage

\begin{consulta}{25}{Copias y alquileres en la tienda 1}
Con dos \codigo{LEFT JOIN}, muestra las películas \codigo{PG-13} de al menos 90
minutos cuyas copias de la tienda 1 sumaron menos de cinco alquileres en mayo de
2005. Incluye las películas con cero alquileres, también si tienen cero copias en
esa tienda. Cuenta cada copia de la tienda 1 una sola vez, se haya alquilado o no
en mayo.
\salida{\codigo{film_id}, \codigo{title}, \codigo{num_ejemplares}, \codigo{num_alquileres}.}{Menor número de alquileres primero. En caso de empate, por título de A a Z.}
\end{consulta}
\resultados{25}{rlrr}
\clearpage

\begin{consulta}{26}{Actores en películas de clasificación R}
Con dos \codigo{RIGHT JOIN}, muestra los actores cuyo apellido empieza por B
y que participaron en al menos cuatro películas de clasificación \codigo{R}.
Cuenta únicamente sus películas de esa clasificación.
\salida{\codigo{actor_id}, \codigo{first_name}, \codigo{last_name}, \codigo{num_peliculas_r}.}{De más a menos películas. Si empatan, por \codigo{actor_id} ascendente.}
\end{consulta}
\resultados{26}{rllr}

\clearpage

\begin{consulta}{27}{Alquileres y pagos de películas R en la tienda 2}
Usa tres \codigo{LEFT JOIN} para enlazar \codigo{film}, \codigo{inventory},
\codigo{rental} y \codigo{payment}.

Muestra las películas de clasificación \codigo{R} que duran al menos 120 minutos
y cuyas copias de la tienda 2 se alquilaron menos de tres veces en mayo de 2005.
Incluye las películas con cero alquileres, incluso si no tienen copias en esa tienda.

Cuenta cada alquiler una sola vez, con independencia de sus pagos. Para esos
alquileres, cuenta y suma únicamente los pagos de importe igual o superior a 4,00.
Si no hay pagos que cumplan ese mínimo, muestra cero en el recuento y en la suma.
\salida{\codigo{film_id}, \codigo{title}, \codigo{num_alquileres},
\codigo{num_pagos}, \codigo{importe_total}.}{De más a menos alquileres. Si empatan, por título de A a Z.}
\end{consulta}
\resultados{27}{rlrrr}

\end{document}
